% sec-07-pharmacology-toxicology.tex - Pharmacology and Toxicology (full stage)
\section{Pharmacology and Toxicology Information}

\subsection{Pharmacology and Drug Distribution}

This section presents the pharmacology and toxicology information supporting the
Phase 1, first-in-human, open-label, single-arm investigation of perioperative
daraxonrasib (RMC-6236) in resectable and borderline-resectable pancreatic
ductal adenocarcinoma (PDAC) described elsewhere in this Investigational New
Drug application (IND v1.0, repository v4.3.0). It is organized and submitted in
accordance with 21 CFR \S 312.23(a)(8), which requires adequate information
about the pharmacological and toxicological studies of the drug, on the basis of
which the sponsor has concluded that it is reasonably safe to conduct the
proposed clinical investigation, including an integrated summary of the
toxicological effects of the drug in animals and in vitro and, where applicable,
a full data tabulation and the names and qualifications of the individuals who
evaluated the studies \cite{phase1guidance}. Consistent with the FDA guidance on
the content and format of an IND for a Phase 1 study, the depth of the
pharmacology and toxicology information presented here is scaled to the early
phase of investigation, to the small number of participants (up to eighteen
treated), and to the short, bounded duration of drug exposure characteristic of
a perioperative dose-finding study \cite{phase1guidance}.

Daraxonrasib is a clinically advanced investigational agent with an established
nonclinical and clinical pharmacology package generated within its late-phase
sponsor program, and it carries FDA Breakthrough Therapy designation granted in
June 2025 for previously treated metastatic PDAC harboring KRAS G12 mutations
\cite{daraxwhipple}. The present IND does not seek to recapitulate that
established package; rather, it cross-references the existing nonclinical
pharmacology and toxicology of daraxonrasib and supplements it with the
study-specific computational pharmacology, drug-distribution, and toxicology
evidence that is unique to the perioperative, curative-intent surgical context
and to the combined drug-device use evaluated here \cite{daraxwhipple,onpremwhippl}.
That supplementary evidence is the principal subject of this section and consists
of three classes of in silico work: an AI digital-twin PDAC simulation
\cite{fdadigtwinpc}, a quantitative systems pharmacology (QSP) metastatic-PDAC
trial simulation with verification, validation, and uncertainty quantification
(VVUQ) \cite{qspmetpancre}, and a 100,000-patient in silico Phase 3 five-arm
PDAC trial conducted in triplicate \cite{chatgpt100kp,pdacdigtwinp}. The combined
regimen also pairs the oral drug with an on-premises large language model (LLM)
directed eight-arm robotic pancreaticoduodenectomy (Whipple) platform that is the
subject of the parallel investigational device exemption content of this
submission; the device has no pharmacologic dose and is not addressed as a
pharmacology or toxicology article here, but the operative and pathology outputs
of the device procedure are inputs to the perioperative drug pause-and-restart
advisory and therefore bear on the drug-distribution discussion that follows
\cite{daraxwhipple}.

The remainder of this section is organized into the three subsections required
by the integrated-summary structure of 21 CFR \S 312.23(a)(8): a pharmacology
summary and conclusions, an integrated toxicology summary, and a full data
tabulation. The pharmacology summary states the RAS(ON) mechanism of action, the
pharmacodynamics, and the drug-distribution basis for the perioperative
pause-and-restart schedule, including the 32-iteration deterministic
pharmacokinetic sweep anchored to a 0.45 ng/mL operative-window baseline serum
concentration. The integrated toxicology summary states the safety margins
supporting the conservative DL1 160 mg once-daily first-in-human starting dose
and the dose-limiting toxicity (DLT) framework that bounds the escalation. The
full data tabulation summarizes the tabulated computational studies that are
provided as supporting documents to this IND. The overall structure of \S 8 and
its relationship to the computational-evidence base is described in Figure 17,
which is rendered with the Investigator Brochure content of \S 5 and is not
reproduced here; the present section relies on tables and prose.

\subsubsection{Pharmacology Summary and Conclusions}

\paragraph{Mechanism of action and pharmacological class.}
Daraxonrasib (RMC-6236) is an orally bioavailable, RAS(ON) multi-selective
(pan-RAS) small-molecule inhibitor. It acts by binding the active, guanosine
triphosphate (GTP) bound state of mutant RAS in a tri-complex with the
ubiquitous chaperone protein cyclophilin A, and the resulting daraxonrasib to
cyclophilin A to RAS(ON) complex sterically occludes the effector-binding
interface, suppressing downstream mitogen-activated protein kinase and
phosphoinositide 3-kinase effector signaling \cite{daraxwhipple}. Because the
agent targets the shared active conformation rather than a single mutant side
chain, it provides broad coverage across the common KRAS oncoprotein variants
found in PDAC, including the G12D, G12V, and G12R alleles that define the
molecular eligibility for this study \cite{daraxwhipple,onpremwhippl}. The
pharmacological class is therefore a RAS(ON) multi-selective inhibitor, and the
mechanistic rationale for evaluating the agent in resectable and
borderline-resectable PDAC is the suppression of micrometastatic outgrowth
during the perioperative window when systemic therapy must otherwise be
interrupted for surgery. The dose and exposure assumptions used throughout this
IND are anchored to the pan-KRAS RASolute 302 program, in which 300 mg once
daily is the established recommended Phase 2 dose, and to the broader late-phase
program comprising RASolute 302 and the first-line RASolve 301 study
\cite{daraxwhipple}. Because daraxonrasib engages the shared GTP-bound
conformation common to the principal mutant alleles rather than discriminating
among individual side chains, the molecular eligibility used in this study, namely
documented KRAS G12 status (G12D, G12V, or G12R), is satisfied by the same broad
pan-KRAS criteria that define eligibility in RASolute 302, and no allele-specific
dose adjustment is anticipated within the studied range
\cite{daraxwhipple,onpremwhippl}. This multi-selectivity is the pharmacologic
property that distinguishes the agent from allele-locked inhibitors and that
motivates its evaluation as the systemic backbone of the perioperative,
curative-intent regimen described in this IND.

\paragraph{Pharmacodynamics.}
The pharmacodynamic profile of daraxonrasib follows directly from its mechanism:
target engagement is a function of the fraction of GTP-bound mutant RAS that is
captured in the inhibitory tri-complex, which in turn tracks systemic drug
exposure. The pharmacodynamic markers relevant to this perioperative study, the
direction of the on-treatment change, and the way each marker is used in the
study are summarized in Table 8. The dominant on-target, class-characteristic
toxicities, principally gastrointestinal events (nausea, diarrhea, stomatitis)
and dermatologic and mucocutaneous events, are themselves pharmacodynamic
consequences of pan-RAS pathway suppression and are managed by the supportive-care
and dose-modification machinery of the protocol rather than by withholding the
mechanism \cite{onpremwhippl}.

\begin{table}[t]
\centering
\caption{Pharmacodynamic markers of daraxonrasib relevant to the perioperative
PDAC study and their use in the protocol.}
\label{tab:sec07-pd}
\begin{tabularx}{\textwidth}{L{3.7cm} L{2.4cm} Y}
\toprule
\textbf{Pharmacodynamic marker} & \textbf{On-treatment direction} & \textbf{Use in this study} \\
\midrule
RAS(ON) tri-complex target engagement & Increases with exposure & Mechanistic basis for exposure-anchored dose escalation; not measured directly in participants but modeled in the QSP and digital-twin evidence \cite{qspmetpancre}. \\
Downstream MAPK / PI3K effector signaling & Suppressed & Mechanistic surrogate for antitumor effect; informs the micrometastatic-outgrowth rationale for restarting drug early postoperatively. \\
Serum daraxonrasib trough concentration & Rises to steady state on continuous dosing; falls during the protocol-mandated operative pause & Objective adherence confirmation and the quantitative input (relative to the 0.5 ng/mL threshold) to the perioperative restart advisory (\S 5, \S 6). \\
Gastrointestinal and dermatologic class effects & Increase with exposure & On-target pharmacodynamic toxicities detected and bounded by the 3+3 DLT framework; managed by protocol supportive care. \\
CA 19-9 tumor marker (exploratory) & Expected to fall with antitumor response & Exploratory biochemical correlate of systemic disease control across the perioperative course. \\
\bottomrule
\end{tabularx}
\end{table}

\figcaption{Table 8. Pharmacodynamic markers of daraxonrasib and their use in the
perioperative PDAC study. Target engagement and effector suppression are
mechanistic; the serum trough is the operative quantitative input to the restart
advisory.}

\paragraph{Drug distribution and the perioperative pharmacokinetic basis.}
Daraxonrasib is administered orally, once daily, in a perioperative schedule
consisting of a defined pre-operative course, a protocol-mandated pause spanning
the operative window, and a postoperative restart whose timing is governed by an
LLM-bound, human-confirmed advisory \cite{onpremwhippl}. The drug-distribution
question that this schedule poses is when, after surgery, systemic exposure
should be re-established so as to suppress micrometastatic outgrowth at the
earliest point that does not jeopardize anastomotic healing. The advisory takes
two inputs, the pancreaticojejunostomy (PJ) International Study Group of
Pancreatic Surgery (ISGPS) postoperative pancreatic-fistula grade (A, B, or C)
and the serum daraxonrasib trough relative to a 0.5 ng/mL threshold, and it
recommends a restart on postoperative day T+7d, T+14d, or T+21d, or a continued
hold \cite{onpremwhippl}. The advisory is bound to a deterministic seed and a
specific repository commit and is always human-confirmed; it is never an
autonomous controller, consistent with the prohibition on full autonomy at 21
CFR \S 312.21(e) \cite{cfr312paiuni}. The drug-distribution logic of the schedule
is therefore explicit and re-traceable: the pre-operative pause is required to
clear systemic exposure ahead of the operation, the operative-window trough is
measured against a single fixed threshold rather than estimated, and the
postoperative restart is keyed jointly to that trough and to the dominant
determinant of fistula risk, the pancreaticojejunostomy, so that systemic
exposure is re-established neither before the highest-risk anastomosis has
declared its healing trajectory nor so late that the therapeutic window for
micrometastatic suppression is lost \cite{onpremwhippl,qspmetpancre}.

The pharmacokinetic basis for this schedule was characterized in a 32-iteration
deterministic simulation sweep, the results of which are summarized in Table 9.
Across all 32 iterations the baseline serum daraxonrasib concentration at the
operative timepoint was 0.45 ng/mL, that is, below the 0.5 ng/mL trough
threshold, reflecting the washout produced by the protocol-mandated
pre-operative pause \cite{onpremwhippl}. Under these sub-threshold operative-window
conditions the advisory recommended a restart on postoperative day T+7d in 29 of
32 iterations (90.6 percent), corresponding to an ISGPS Grade A fistula with a
sub-threshold trough; it recommended T+14d in 3 of 32 iterations (9.4 percent),
corresponding to an ISGPS Grade B fistula or a delayed serum rise; and it
recommended T+21d in 0 of 32 iterations (0.0 percent) in this sweep. An ISGPS
Grade C fistula maps to a T+21d restart or a continued hold and to the drug
stop rule of the protocol. The restart-day distribution and the uniform 0.45
ng/mL baseline are the quantitative drug-distribution conclusions that support
the perioperative schedule and the advisory thresholds.

\begin{table}[t]
\centering
\caption{Perioperative pharmacokinetic sweep: 32-iteration deterministic
restart-advisory distribution at a uniform 0.45 ng/mL operative-window baseline
serum daraxonrasib concentration.}
\label{tab:sec07-pk}
\begin{tabularx}{\textwidth}{L{2.5cm} L{1.9cm} L{2.0cm} Y}
\toprule
\textbf{Advisory output} & \textbf{Iterations (of 32)} & \textbf{Proportion} & \textbf{Determining condition} \\
\midrule
Restart T+7d & 29 & 90.6\% & ISGPS Grade A fistula with sub-threshold trough (serum 0.45 ng/mL $<$ 0.5 ng/mL); earliest safe re-establishment of systemic exposure. \\
Restart T+14d & 3 & 9.4\% & ISGPS Grade B fistula or a delayed serum rise; restart deferred one week to protect anastomotic healing. \\
Restart T+21d & 0 & 0.0\% & ISGPS Grade C fistula maps here or to a continued hold; not triggered in this sweep given the uniform Grade A / B inputs. \\
Continued hold & 0 & 0.0\% & Grade C fistula with persistent risk; defers to the protocol drug stop rule. \\
\midrule
\textbf{Baseline (all 32)} & 32 & 100\% & Operative-window serum daraxonrasib 0.45 ng/mL (sub-threshold) in every iteration, reflecting pre-operative washout. \\
\bottomrule
\end{tabularx}
\end{table}

\figcaption{Table 9. The 32-iteration deterministic perioperative pharmacokinetic
sweep. The operative-window baseline serum daraxonrasib concentration was 0.45
ng/mL (below the 0.5 ng/mL trough threshold) in all iterations; the advisory
recommended T+7d in 29 of 32, T+14d in 3 of 32, and T+21d in 0 of 32.}

\paragraph{Pharmacology conclusions.}
On the basis of the established RAS(ON) multi-selective mechanism, the broad
pan-KRAS coverage encompassing the G12 alleles that predominate in PDAC, the
exposure-anchored pharmacodynamics summarized in Table 8, and the deterministic
perioperative pharmacokinetic sweep summarized in Table 9, the sponsor concludes
that the pharmacology of daraxonrasib supports the proposed perioperative,
once-daily oral schedule and the exposure thresholds used by the restart
advisory. The 0.45 ng/mL sub-threshold operative-window baseline and the
resulting predominant T+7d restart recommendation provide a quantitative basis
for re-establishing systemic exposure early, while the deferral logic for higher
fistula grades preserves anastomotic healing, so that the schedule is reasonably
expected to balance micrometastatic suppression against surgical risk in the
study population \cite{onpremwhippl,qspmetpancre}.

\subsubsection{Toxicology: Integrated Summary}

\paragraph{Conservative first-in-human starting dose and its safety margin.}
The first-in-human starting dose in this study is DL1, 160 mg once daily, which
is deliberately set below the RASolute 302 reference recommended Phase 2 dose of
300 mg once daily \cite{daraxwhipple}. The escalation then steps up through DL2
at 220 mg once daily to DL3 at the 300 mg once-daily reference dose, so that the
entire dosing range evaluated in this study (160 mg to 300 mg) lies at or below
an exposure that has already been characterized in the late-phase sponsor
program \cite{daraxwhipple,onpremwhippl}. The starting dose therefore carries a
built-in safety margin in two senses: it is approximately 53 percent of the
established reference dose by milligram amount (160 of 300 mg), and it sits
within, rather than above, the exposure range for which the established
nonclinical and clinical toxicology of daraxonrasib already supports
administration. Because the agent is clinically advanced and its established
toxicology package supports dosing at and above the doses studied here, the
conservative DL1 starting dose, combined with the bounded perioperative exposure
duration, provides the integrated-summary basis for the sponsor's conclusion that
it is reasonably safe to begin the proposed clinical investigation
\cite{phase1guidance,daraxwhipple}.

\paragraph{Dose-limiting toxicity framework.}
The toxicology of the escalation is bounded prospectively by a standard 3+3
dose-limiting toxicity (DLT) framework with sentinel staggered enrollment, rather
than retrospectively. Three participants are enrolled at each dose level and
observed across a 28-day DLT window. If no DLT occurs among the first three, the
cohort escalates; if one of three experiences a DLT, the cohort expands to six; a
single DLT among six permits escalation, whereas two or more DLTs at a level fix
the maximum tolerated dose (MTD) at the preceding level. The MTD is the highest
level at which fewer than two of six DLT-evaluable participants experience a DLT,
and the recommended Phase 2 dose is selected at or below the MTD
\cite{onpremwhippl}. The dose levels and their relationship to the reference
dose and the DLT framework are summarized in Table 10. The class-characteristic
on-target toxicities of pan-RAS inhibition, principally gastrointestinal events
(nausea, diarrhea, stomatitis) and dermatologic and mucocutaneous events, are
the toxicities most likely to manifest as DLTs and are precisely the events the
28-day window is designed to detect and bound; perioperative exposure
additionally raises a theoretical risk to wound and anastomotic healing if
restart timing is suboptimal, which is the specific hazard the pause-and-restart
advisory (Table 9) is designed to mitigate \cite{onpremwhippl}.

\begin{table}[t]
\centering
\caption{Integrated toxicology summary: dose levels, safety margin relative to
the reference dose, and the 3+3 dose-limiting toxicity framework.}
\label{tab:sec07-tox}
\begin{tabularx}{\textwidth}{L{1.4cm} L{2.1cm} L{2.4cm} Y}
\toprule
\textbf{Level} & \textbf{Dose (oral QD)} & \textbf{Fraction of 300 mg reference} & \textbf{Toxicology basis and escalation rule} \\
\midrule
DL1 & 160 mg & 53\% & Conservative first-in-human starting dose; below the RASolute 302 reference; 3 enrolled, 28-day DLT window; 0/3 escalate, 1/3 expand to 6. \\
DL2 & 220 mg & 73\% & Intermediate level; entered only after DL1 clears the DLT window; same 3+3 expansion and escalation logic. \\
DL3 & 300 mg & 100\% & Equal to the established RASolute 302 recommended Phase 2 dose; highest level studied; MTD if fewer than 2/6 DLTs. \\
\midrule
\multicolumn{4}{>{\raggedright\arraybackslash}p{\dimexpr\textwidth-2\tabcolsep}}{\textbf{DLT rule.} MTD is the highest level with fewer than 2 of 6 DLT-evaluable DLTs; RP2D is selected at or below the MTD. Two device-or-procedure related deaths at any time, or one device-or-procedure serious adverse event within a cohort, trigger DSMB review and a binding halt.} \\
\bottomrule
\end{tabularx}
\end{table}

\figcaption{Table 10. Integrated toxicology summary. The 160 mg starting dose is
53 percent of the 300 mg reference dose, and the entire 160 mg to 300 mg range
lies at or below an exposure characterized in the established daraxonrasib
program; the 3+3 framework bounds the escalation prospectively.}

\paragraph{Computational toxicology and QSP evidence.}
The integrated toxicology conclusion is reinforced by study-specific computational
toxicology evidence. A quantitative systems pharmacology (QSP) metastatic-PDAC
trial simulation, conducted with code, verification, validation, and uncertainty
quantification (VVUQ), and an accompanying playbook, models the exposure to
response and exposure to toxicity relationships of daraxonrasib in a mechanistic
systems framework and supports the plausibility of the exposure thresholds and
the dose range studied here \cite{qspmetpancre}. An AI digital-twin
pancreatic-cancer simulation, developed for in silico regulatory-compliance and
cost-efficiency purposes, provides a complementary patient-level
toxicodynamic representation \cite{fdadigtwinpc}. These computational works do
not substitute for the established nonclinical toxicology of daraxonrasib;
rather, they supplement it with study-specific, in silico characterization of the
perioperative context, and they are tabulated as supporting documents in the full
data tabulation below \cite{qspmetpancre,fdadigtwinpc}.

\paragraph{Integrated toxicology conclusion.}
Taking together the conservative 160 mg first-in-human starting dose set at 53
percent of an already-characterized reference exposure, the prospective 3+3 DLT
framework with a 28-day window and sentinel staggered enrollment, the binding
halt rules, the perioperative pause-and-restart advisory that mitigates the
theoretical anastomotic-healing hazard, and the supporting QSP and digital-twin
computational toxicology, the sponsor concludes that the toxicological profile of
daraxonrasib supports the conclusion that it is reasonably safe to conduct the
proposed Phase 1 investigation at the doses and schedule specified, within the
meaning of 21 CFR \S 312.23(a)(8) \cite{phase1guidance,cfr312paiuni}.

\subsubsection{Toxicology: Full Data Tabulation}

\paragraph{Scope and form of the tabulation.}
For a clinically advanced agent such as daraxonrasib, the full data tabulation
required by 21 CFR \S 312.23(a)(8)(ii) for the established nonclinical
toxicology studies is provided by cross-reference to the sponsor's existing
nonclinical package, and the qualifications of the individuals who evaluated
those studies are stated in that package \cite{phase1guidance,daraxwhipple}.
What is tabulated here, and provided as supporting documents to this IND, is the
study-specific computational-evidence base that characterizes daraxonrasib and
the perioperative drug-device regimen in silico. These works were generated by
the sponsor-investigator over 2025 and 2026 and comprise the digital-twin, QSP,
and large-scale in silico Phase 3 simulations described above. Their identity,
the evidence each contributes, and the location of each as a supporting document
are tabulated in Table 11.

\begin{table}[t]
\centering
\caption{Full data tabulation: study-specific computational-evidence supporting
documents for the pharmacology and toxicology of the perioperative daraxonrasib
PDAC regimen.}
\label{tab:sec07-comp}
\begin{tabularx}{\textwidth}{L{3.6cm} L{1.9cm} Y}
\toprule
\textbf{Computational study} & \textbf{Form} & \textbf{Evidence contributed to \S 8} \\
\midrule
AI digital-twin PDAC simulation \cite{fdadigtwinpc} & Patient-level digital twin & In silico, regulatory-compliance-oriented toxicodynamic and disease-progression representation of PDAC; supports the perioperative drug-distribution and toxicology reasoning at the individual-patient level. \\
QSP metastatic-PDAC trial simulation with VVUQ \cite{qspmetpancre} & Mechanistic QSP model + VVUQ + playbook & Exposure-to-response and exposure-to-toxicity relationships in a verified, validated, uncertainty-quantified systems-pharmacology framework; underpins the dose range and exposure thresholds. \\
100,000-patient in silico Phase 3, 5-arm PDAC, triplicate \cite{chatgpt100kp,pdacdigtwinp} & Large-scale virtual trial (triplicate) & Population-scale in silico characterization of a five-arm PDAC Phase 3 across 100,000 virtual patients, run in triplicate; supports the population-level plausibility of the regimen and the trial design. \\
End-to-end PDAC digital-twin trial proposals \cite{pdacdigtwinp} & Trial-design proposals & End-to-end digital-twin clinical-trial scaffolding adapted for the perioperative drug-device design described in this IND. \\
\midrule
\multicolumn{3}{>{\raggedright\arraybackslash}p{\dimexpr\textwidth-2\tabcolsep}}{\textbf{Cross-reference.} The established nonclinical toxicology full data tabulation for daraxonrasib (RMC-6236) is provided by reference to the sponsor's existing nonclinical package; the qualifications of the evaluators are stated there. The 32-iteration perioperative pharmacokinetic sweep (Table 9) is tabulated above.} \\
\bottomrule
\end{tabularx}
\end{table}

\figcaption{Table 11. Full data tabulation of the study-specific computational
evidence (digital-twin, QSP with VVUQ, and the 100,000-patient in silico Phase 3
PDAC triplicate) provided as supporting documents, with the established
nonclinical toxicology incorporated by cross-reference.}

\paragraph{Evaluators and incorporation by reference.}
The study-specific computational studies tabulated in Table 11 were designed,
executed, and evaluated by the sponsor-investigator and the study computational
team, whose qualifications are stated in the Investigator and Facilities data of
this submission; the verification, validation, and uncertainty-quantification
methodology applied to the QSP work is documented in that study's accompanying
VVUQ record \cite{qspmetpancre}. The established nonclinical pharmacology and
toxicology of daraxonrasib, including the full data tabulations and evaluator
qualifications for those studies, is incorporated by reference to the sponsor's
existing nonclinical package consistent with the Phase 1 IND content guidance
\cite{phase1guidance,daraxwhipple}. As the investigation proceeds, any new
nonclinical or computational toxicology information bearing on participant safety
will be submitted through information amendments under 21 CFR \S 312.31 and, where
a safety signal meets the reporting thresholds, through IND safety reports under
21 CFR \S 312.32 \cite{cfr312paiuni}.

\paragraph{Conclusion of the pharmacology and toxicology information.}
The pharmacology summary (\S 8.1.1), the integrated toxicology summary
(\S 8.1.2), and the full data tabulation (\S 8.1.3) together provide adequate
information about the pharmacological and toxicological properties of
daraxonrasib, scaled to a Phase 1 study, on the basis of which the sponsor has
concluded that it is reasonably safe to conduct the proposed perioperative,
once-daily oral investigation in resectable and borderline-resectable PDAC at the
160 mg, 220 mg, and 300 mg dose levels under the 3+3 dose-limiting toxicity
framework, in satisfaction of 21 CFR \S 312.23(a)(8)
\cite{phase1guidance,cfr312paiuni,daraxwhipple}.
